\chapter{Modèles probabilistes élémentaires}
\label{workbook:chapter:08}
\section{Expérience, univers et événement}
\label{workbook:section:08.01}
\begin{exerciseblock}{Expérience, univers et événement}
\item Pour deux dés, définir un univers équiprobable et l'événement $A=$ « somme égale à $8$ ».
\item Avec une pièce lancée trois fois, écrire les événements $B=$ « au moins deux faces » et $C=$ « aucun pile ».
\item Exprimer en mots $A\cap B$, $A\cup B$ et $A^c$ dans un contexte de cartes où $A=$ « rouge » et $B=$ « figure ».
\end{exerciseblock}
\section{Axiomes et équiprobabilité}
\label{workbook:section:08.02}
\begin{exerciseblock}{Axiomes et équiprobabilité}
\item Vérifier que $P(\Omega)=1$, $P(\varnothing)=0$ et que $P(A^c)=1-P(A)$ sont cohérents pour un dé équilibré.
\item Calculer la probabilité d'obtenir une somme au moins $10$ avec deux dés.
\item Donner un exemple où les résultats listés ne sont pas équiprobables et où la formule $|A|/|\Omega|$ serait incorrecte.
\end{exerciseblock}
\section{Probabilité conditionnelle}
\label{workbook:section:08.03}
\begin{exerciseblock}{Probabilité conditionnelle}
\item Dans un jeu de $52$ cartes, calculer $P(\text{roi}\mid\text{figure})$.
\item Une classe compte $18$ femmes dont $6$ étudient les maths, et $12$ hommes dont $3$ étudient les maths. Calculer $P(M\mid F)$ et $P(F\mid M)$.
\item Réécrire la formule $P(A\mid B)=P(A\cap B)/P(B)$ comme une formule du produit.
\end{exerciseblock}
\section{Indépendance}
\label{workbook:section:08.04}
\begin{exerciseblock}{Indépendance}
\item Tester l'indépendance de $A=$ « premier dé pair » et $B=$ « somme des deux dés paire ».
\item Expliquer pourquoi deux événements incompatibles de probabilités positives ne peuvent pas être indépendants.
\item Une pièce biaisée a $P(F)=0{,}6$. Pour deux lancers indépendants, calculer la probabilité d'obtenir exactement une face.
\end{exerciseblock}
\section{Épreuve de Bernoulli et modèle binomial}
\label{workbook:section:08.05}
\begin{exerciseblock}{Épreuve de Bernoulli et modèle binomial}
\item Identifier $n$, $p$ et la variable aléatoire dans une suite de $12$ tirs avec probabilité de réussite $0{,}7$.
\item Calculer la probabilité d'exactement $8$ réussites.
\item Calculer la probabilité d'au moins une réussite en $5$ essais lorsque $p=0{,}2$, par complément.
\end{exerciseblock}
\section{Espérance mathématique}
\label{workbook:section:08.06}
\begin{exerciseblock}{Espérance mathématique}
\item Une variable $X$ prend $0,1,2$ avec probabilités $0{,}2,0{,}5,0{,}3$. Calculer $E(X)$.
\item Un jeu coûte $4\,$\$; il rapporte $20\,$\$ avec probabilité $0{,}1$ et rien sinon. Calculer le gain net espéré.
\item Pour $X\sim\mathrm{Bin}(n,p)$, interpréter $E(X)=np$ sans le confondre avec un résultat certain.
\end{exerciseblock}
\section{Proportion, fréquence et probabilité}
\label{workbook:section:08.07}
\begin{exerciseblock}{Proportion, fréquence et probabilité}
\item Une fréquence de succès vaut $63/100$ après $100$ essais puis $617/1000$ après $1000$ essais. Comparer les deux estimations de $p$.
\item Expliquer la différence entre fréquence observée et probabilité du modèle.
\item Simuler conceptuellement ce qui arrive à la variabilité de la fréquence lorsque le nombre d'essais augmente.
\end{exerciseblock}
\section{Le conditionnement avec des effectifs concrets}
\label{workbook:section:08.08}
\begin{exerciseblock}{Le conditionnement avec des effectifs concrets}
\item Dans $200$ dossiers, $30$ ont un test positif; parmi eux $24$ sont réellement malades. Parmi les $170$ négatifs, $6$ sont malades. Construire le tableau à double entrée.
\item Calculer $P(M\mid +)$, $P(+\mid M)$ et $P(M)$.
\item Expliquer pourquoi confondre $P(M\mid +)$ avec $P(+\mid M)$ peut conduire à une interprétation erronée.
\end{exerciseblock}
\section{Interpréter une information probabiliste}
\label{workbook:section:08.09}
\begin{exerciseblock}{Interpréter une information probabiliste}
\item Un risque passe de $2\%$ à $1\%$. Calculer réduction absolue et réduction relative.
\item Une étude annonce « deux fois plus probable » sans donner le risque de base. Expliquer ce qui manque.
\item Formuler trois questions à poser avant d'accepter une affirmation probabiliste issue d'un sondage.
\end{exerciseblock}
\section*{Synthèse du chapitre}
\begin{synthesisbox}
\item Une maladie touche $1\%$ d'une population. Un test a une sensibilité de $95\%$ et un taux de faux positifs de $5\%$. Sur $10\,000$ personnes, construire les effectifs et calculer $P(M\mid+)$.
\item On lance un dé jusqu'à obtenir un six, au maximum trois fois. Construire l'arbre, calculer la probabilité de succès et le nombre moyen de lancers effectués.
\item Un jeu donne $10\,$\$ avec probabilité $0{,}2$, $3\,$\$ avec probabilité $0{,}5$, et $0$ sinon. Quel prix d'entrée rend le jeu équitable?
\end{synthesisbox}
